\paragrafos{Πολλαπλασιασμός ακεραίων: το μείον επί μείον}{Α.2.4, «Πολλαπλασιασμός ακεραίων»}

\begin{anagke}
Το βράδυ στο Μέτσοβο η θερμοκρασία πέφτει $2$ βαθμούς κάθε ώρα. Η Πολυξένη παρατηρεί το θερμόμετρο στις $20{:}00$, που δείχνει $0$°. «Τι θα δείχνει σε τρεις ώρες;». Ο Νίκος σκέφτεται: «Κάθε ώρα $-2$, τρεις ώρες: $-2-2-2=-6$». Η Αγγελική το γράφει ως πολλαπλασιασμό: $3\cdot(-2)=-6$. «Και πριν από τρεις ώρες, στις $17{:}00$;» ρωτάει η Ρήνη. «Αφού τώρα πέφτει δύο βαθμούς ανά ώρα, πριν από τρεις ώρες είχε $6$ βαθμούς περισσότερο», απαντάει η Πολυξένη. Ο ρυθμός είναι $-2$ και ο χρόνος «τρεις ώρες πριν» είναι $-3$, και το αποτέλεσμα είναι $+6$. Άρα $(-2)\cdot(-3)=+6$.

Ο Στέλιος ξαφνιάζεται: «Το μείον επί μείον κάνει συν; Σαν τη διπλή άρνηση της γλώσσας που καταλήγει σε κατάφαση;». Η κυρία Αλεξάνδρα του απαντά: «Η γλώσσα είναι ωραία αναλογία, αλλά τα μαθηματικά θέλουν να δουν γιατί». Το ίδιο το μοτίβο δείχνει το γιατί: $3\cdot(-2)=-6$, $2\cdot(-2)=-4$, $1\cdot(-2)=-2$, $0\cdot(-2)=0$. Κάθε φορά που το πρώτο μικραίνει κατά $1$ το γινόμενο μεγαλώνει κατά $2$. Συνεχίζοντας: $(-1)\cdot(-2)=2$, $(-2)\cdot(-2)=4$.

Σε αυτή την παράγραφο η παρέα μαθαίνει τον κανόνα των προσήμων στον πολλαπλασιασμό, τον εφαρμόζει σε γινόμενα με πολλούς παράγοντες και στις παραστάσεις με γράμματα, και προσέχει να μη μπερδεύει τον κανόνα του πολλαπλασιασμού με αυτόν της πρόσθεσης.
\end{anagke}

\begin{lexi}{παράγοντας, γινόμενο}{παράγοντας, γινόμενο}
\lexiword{Παράγοντες} είναι οι αριθμοί που πολλαπλασιάζονται και \lexiword{γινόμενο} αυτό που «γίνεται». Ο \lexiword{πολλαπλασιασμός} έρχεται από το «πολλαπλάσιο» και το «πλάσσω»: φτιάχνω πολλές φορές το ίδιο.
\end{lexi}

\begin{metafrasi}
Πέφτει $2$ βαθμούς την ώρα, για $3$ ώρες. & $3\cdot(-2)=-6$\\
Πριν από $3$ ώρες. & $(-2)\cdot(-3)=+6$\\
Ομόσημοι παράγοντες. & γινόμενο θετικό\\
Ετερόσημοι παράγοντες. & γινόμενο αρνητικό\\
Ένας παράγοντας μηδέν. & γινόμενο $0$\\
Το γινόμενο των $\alpha$ και $\beta$. & $\alpha\beta$ ή $\alpha\cdot\beta$\\
\end{metafrasi}

\tmima{Η θεωρία}

\tmimasmall{Γινόμενο δύο ακεραίων}

Σε ένα παιχνίδι γνώσεων, ο παίκτης κερδίζει $5$ πόντους ($+5$) για κάθε σωστή απάντηση και χάνει $2$ πόντους ($-2$) για κάθε λάθος. Αν απαντήσει σωστά σε $3$ ερωτήσεις έχει $(+3)\cdot(+5)=(+5)+(+5)+(+5)=+15$ πόντους. Αν απαντήσει λάθος σε $3$ ερωτήσεις έχει $(+3)\cdot(-2)=(-2)+(-2)+(-2)=-6$ πόντους. Παρατηρούμε ότι:

\begin{itemize}[leftmargin=1.2em,itemsep=0.3ex,topsep=2pt]
\item Το γινόμενο δύο θετικών είναι θετικός αριθμός.
\item Το γινόμενο ενός θετικού και ενός αρνητικού είναι αρνητικός αριθμός.
\end{itemize}

Για τους δύο αρνητικούς, ας δούμε το μοτίβο:
$$\begin{array}{rcl}
(-4)\cdot(+3)&=&-12\\
(-4)\cdot(+2)&=&-8\\
(-4)\cdot(+1)&=&-4\\
(-4)\cdot\phantom{(}0&=&0\\
(-4)\cdot(-1)&=&\ \ldots\\
(-4)\cdot(-2)&=&\ \ldots
\end{array}$$
Κάθε γραμμή δίνει αποτέλεσμα κατά $4$ μεγαλύτερο από την προηγούμενη. Αν συνεχίσουμε με τον ίδιο τρόπο, πρέπει να είναι $(-4)\cdot(-1)=+4$ και $(-4)\cdot(-2)=+8$. Διαπιστώνουμε ότι το γινόμενο δύο αρνητικών είναι θετικός αριθμός.

\begin{kanonas}{γινόμενο ακεραίων}
Για να πολλαπλασιάσω δύο ακέραιους, πολλαπλασιάζω τις απόλυτες τιμές τους και στο γινόμενο βάζω:
\begin{itemize}[leftmargin=1.2em,itemsep=0.3ex,topsep=2pt]
\item πρόσημο $+$, αν είναι \textbf{ομόσημοι}: $(+2)\cdot(+3)=+6$ και $(-2)\cdot(-3)=+6$.
\item πρόσημο $-$, αν είναι \textbf{ετερόσημοι}: $(+4)\cdot(-5)=-20$ και $(-3)\cdot(+3)=-9$.
\end{itemize}
Αν ένας παράγοντας είναι $0$, το γινόμενο είναι $0$.
\end{kanonas}

Σύντομα: $(+)\cdot(+)=+$, $(+)\cdot(-)=-$, $(-)\cdot(+)=-$, $(-)\cdot(-)=+$. Το $+$ δίνεται όταν τα πρόσημα είναι ίδια. Όταν χρησιμοποιούμε γράμματα, γράφουμε το γινόμενο χωρίς τελεία: $\alpha\beta$ αντί $\alpha\cdot\beta$ και $3\gamma$ αντί $3\cdot\gamma$.

\begin{center}
\begin{tikzpicture}[font=\small]
  \draw[rounded corners=3pt,fill=gblueL,draw=gblue] (0,0) rectangle (9.6,3.2);
  \node[anchor=west] at (0.3,2.65) {Ο φίλος του φίλου μου είναι φίλος μου.};
  \node[anchor=east] at (9.3,2.65) {$+\cdot+=+$};
  \node[anchor=west] at (0.3,1.95) {Ο εχθρός του φίλου μου είναι εχθρός μου.};
  \node[anchor=east] at (9.3,1.95) {$+\cdot-=-$};
  \node[anchor=west] at (0.3,1.25) {Ο φίλος του εχθρού μου είναι εχθρός μου.};
  \node[anchor=east] at (9.3,1.25) {$-\cdot+=-$};
  \node[anchor=west] at (0.3,0.55) {Ο εχθρός του εχθρού μου είναι φίλος μου.};
  \node[anchor=east] at (9.3,0.55) {$-\cdot-=+$};
\end{tikzpicture}
\end{center}

\tmimasmall{Γινόμενο πολλών παραγόντων}

Για να υπολογίσω ένα γινόμενο πολλών παραγόντων, πολλαπλασιάζω τις απόλυτες τιμές τους και στο γινόμενο βάζω:
\begin{itemize}[leftmargin=1.2em,itemsep=0.3ex,topsep=2pt]
\item πρόσημο $+$, αν το πλήθος των \textbf{αρνητικών} παραγόντων είναι \textbf{άρτιο},
\item πρόσημο $-$, αν το πλήθος των \textbf{αρνητικών} παραγόντων είναι \textbf{περιττό}.
\end{itemize}
Το πρόσημο λοιπόν το καθορίζει μόνο το πλήθος των αρνητικών παραγόντων. Αν τουλάχιστον ένας παράγοντας είναι $0$, το γινόμενο είναι $0$. Για παράδειγμα $(-2)\cdot3\cdot(-1)\cdot(-5)\cdot2=-60$ (τρεις αρνητικοί) και $(-1)\cdot(-1)\cdot(-1)\cdot(-1)=+1$ (τέσσερις αρνητικοί).

Ισχύουν οι γνωστές ιδιότητες: $\alpha\beta=\beta\alpha$, $(\alpha\beta)\gamma=\alpha(\beta\gamma)$, $\alpha\cdot1=\alpha$, $\alpha\cdot0=0$ και $\alpha\cdot(-1)=-\alpha$.

\begin{prosoxi}
\begin{itemize}[leftmargin=1.2em,itemsep=0.3ex,topsep=2pt]
\item Στον πολλαπλασιασμό δύο αρνητικοί δίνουν θετικό, ενώ στην πρόσθεση δύο αρνητικοί δίνουν αρνητικό: $(-3)\cdot(-4)=+12$ αλλά $(-3)+(-4)=-7$.
\item Το πρόσημο του γινομένου μετράει το πλήθος των \emph{αρνητικών} παραγόντων, όχι το πλήθος των θετικών.
\item Ο πολλαπλασιασμός δεν «μεγαλώνει» πάντα: $5\cdot(-2)=-10<5$.
\item Στις παραστάσεις με πράξεις, ο πολλαπλασιασμός προηγείται της πρόσθεσης και της αφαίρεσης: $5-3\cdot(-2)=5-(-6)=11$, όχι $2\cdot(-2)=-4$.
\item Αν στο γινόμενο υπάρχει παράγοντας $0$, δεν χρειάζεται να μετρήσω πρόσημα: το γινόμενο είναι $0$.
\end{itemize}
\end{prosoxi}

\tmima{Λυμένα παραδείγματα}

\paradeigma{Γινόμενα δύο παραγόντων}
\tixero Τα γινόμενα $(+6)\cdot(+7)$, $(+4)\cdot(-5)$, $(-8)\cdot(+3)$, $(-6)\cdot(-7)$, $(-15)\cdot0$.\par
\tithelo Να τα υπολογίσω.\par
\possindeo Πολλαπλασιάζω τις απόλυτες τιμές και βάζω $+$ αν τα πρόσημα είναι ίδια, $-$ αν είναι διαφορετικά.\par
\apantisi $42$, $-20$, $-24$, $42$, $0$.\par
\elegxos Στα ετερόσημα ($-20$ και $-24$) το αποτέλεσμα είναι αρνητικό και στα ομόσημα θετικό.

\paradeigma{Πολλοί παράγοντες}
\tixero Τα γινόμενα $(-2)\cdot3\cdot(-1)\cdot(-5)\cdot2$ και $2\cdot(-3)\cdot0\cdot(-7)$.\par
\tithelo Να τα υπολογίσω.\par
\possindeo Στο πρώτο μετράω τους αρνητικούς: $-2$, $-1$, $-5$, δηλαδή τρεις (περιττό πλήθος), άρα το γινόμενο είναι αρνητικό. Η απόλυτη τιμή είναι $2\cdot3\cdot1\cdot5\cdot2=60$. Το δεύτερο έχει παράγοντα $0$.\par
\apantisi $-60$ και $0$.\par
\elegxos Αν άλλαζα έναν παράγοντα σε θετικό, θα είχα δύο αρνητικούς και το αποτέλεσμα θα γινόταν θετικό.

\paradeigma{Πόντοι στο παιχνίδι}
\tixero Κάθε πράσινη κάρτα μετράει $+3$ πόντους και κάθε μαύρη $-2$. Η Φανή έχει $4$ πράσινες και $3$ μαύρες.\par
\tithelo Πόσους πόντους έχει;\par
\possindeo $4\cdot(+3)+3\cdot(-2)=12+(-6)=6$.\par
\apantisi $6$ πόντους.\par
\elegxos Οι πράσινες δίνουν $12$ και οι μαύρες αφαιρούν $6$: $12-6=6$.

\paradeigma{Παράσταση με πράξεις και παρενθέσεις}
\tixero Η παράσταση $5\cdot(-2+6)-3\cdot(-7+4)$.\par
\tithelo Να την υπολογίσω.\par
\possindeo Πρώτα οι παρενθέσεις: $-2+6=4$ και $-7+4=-3$. Μετά οι πολλαπλασιασμοί: $5\cdot4=20$ και $3\cdot(-3)=-9$. Τέλος η αφαίρεση: $20-(-9)=29$.\par
\apantisi $29$.\par
\elegxos Το δεύτερο γινόμενο είναι αρνητικό και το αφαιρούμε: ένα «μείον» μπροστά σε αρνητικό δίνει $+$.

\paradeigma{Αριθμητικές τιμές παραστάσεων}
\tixero Οι τιμές $\alpha=-2$, $\beta=+5$, $\gamma=-3$.\par
\tithelo Να βρω τις τιμές των $\alpha\beta$, $\alpha\gamma$, $\alpha\beta\gamma$ και $3\beta-\alpha\gamma$.\par
\possindeo $\alpha\beta=(-2)\cdot5=-10$. $\alpha\gamma=(-2)\cdot(-3)=6$. $\alpha\beta\gamma=(-10)\cdot(-3)=30$. $3\beta-\alpha\gamma=3\cdot5-6=15-6=9$.\par
\apantisi $-10$, $6$, $30$, $9$.\par
\elegxos Στο $\alpha\beta\gamma$ οι αρνητικοί παράγοντες είναι δύο (άρτιο πλήθος): θετικό γινόμενο.

\paradeigma{Ποιοι δύο ακέραιοι;}
\tixero Δύο ακέραιοι έχουν γινόμενο $+12$ και άθροισμα $-7$.\par
\tithelo Να τους βρω.\par
\possindeo Αφού το γινόμενο είναι θετικό, είναι ομόσημοι. Αφού το άθροισμα είναι αρνητικό, είναι και οι δύο αρνητικοί. Ζεύγη με γινόμενο $12$: $(-1,\,-12)$, $(-2,\,-6)$, $(-3,\,-4)$. Το άθροισμα $-7$ το δίνει το τελευταίο.\par
\apantisi $-3$ και $-4$.\par
\elegxos $(-3)\cdot(-4)=12$ και $(-3)+(-4)=-7$.

\tmima{Ασκήσεις}

\zoni{Το κατάλαβα}{κανόνας προσήμων, πολλοί παράγοντες}

\ex Υπολόγισε τα γινόμενα: α) $(+2)\cdot(+3)$ \quad β) $(-4)\cdot(+6)$ \quad γ) $(-6)\cdot(-7)$ \quad δ) $(+5)\cdot(-5)$ \quad ε) $(-15)\cdot0$ \quad στ) $(-1)\cdot(+207)$.

\ex Συμπλήρωσε με τη σωστή λέξη. α) Το γινόμενο δύο ομοσήμων ακεραίων είναι πάντα \ldots\ αριθμός. β) Το γινόμενο δύο ετεροσήμων ακεραίων είναι πάντα \ldots\ αριθμός. γ) Το πρόσημο ενός γινομένου καθορίζεται από το πλήθος των \ldots\ παραγόντων. δ) Αν το γινόμενο δύο ακεραίων είναι αρνητικό, οι αριθμοί είναι \ldots.

\ex Υπολόγισε: α) $(-1)\cdot(-1)\cdot(-1)\cdot(-1)\cdot(-1)$ \quad β) $(-1)\cdot(-2)\cdot(+3)\cdot(-4)\cdot(-5)$ \quad γ) $(+5)\cdot(-10)\cdot(+5)\cdot(-10)\cdot(-2)$.

\ex Συμπλήρωσε τον πίνακα, όπου κάθε κελί είναι το γινόμενο της γραμμής με τη στήλη.
\begin{center}
\renewcommand{\arraystretch}{1.3}
\begin{tabular}{c|c|c|c|c}
$\cdot$ & $+3$ & $-4$ & $-6$ & $0$\\\hline
$-5$ & & & &\\\hline
$-1$ & & & &\\\hline
$+2$ & & & &\\\hline
$7$ & & & &\\
\end{tabular}
\end{center}

\ex Βρες τον αριθμό που λείπει ώστε να αληθεύουν οι ισότητες: α) $(+3)\cdot(-4)=(-2)\cdot(\ldots)$ \quad β) $(-5)\cdot(-6)=(+3)\cdot(\ldots)$ \quad γ) $(+9)\cdot(+8)=(-6)\cdot(\ldots)$ \quad δ) $(-3)\cdot8=12\cdot(\ldots)$.

\ex Υπολόγισε: α) $2\cdot(-5+3)+4\cdot(-1+6)$ \quad β) $(-7+11-13)\cdot(-5+4)$ \quad γ) $(+6)\cdot(-3)+(-1)\cdot(+4)-2\cdot(-5)\cdot(-4)$ \quad δ) $3\cdot\big[\,2\cdot(-4)-(-5)\cdot2\,\big]$.

\ex Σωστό ή λάθος; α) Το γινόμενο ενός θετικού και ενός αρνητικού είναι αρνητικό. \quad β) $(-3)\cdot(-4)=-12$. \quad γ) Αν ένας παράγοντας είναι $0$, το γινόμενο είναι $0$. \quad δ) Το γινόμενο $(-1)\cdot(-1)\cdot\ldots\cdot(-1)$ με $100$ παράγοντες είναι $-1$. \quad ε) $3\cdot(-2)=(-2)\cdot3$. \quad στ) Το γινόμενο αρνητικών ακεραίων είναι πάντα αρνητικό. \quad ζ) $\alpha\cdot(-1)=-\alpha$.

\ex Για $\alpha=-1$, $\beta=+4$, $\gamma=-5$ υπολόγισε τα $\alpha\beta$, $\beta\gamma$, $\alpha\beta\gamma$, $\alpha\gamma-\beta$ και $\alpha-\beta\gamma$.

\zoni{Το χρησιμοποιώ}{θερμόμετρα, ασανσέρ, παιχνίδια}

\ex \textbf{Η νύχτα.} Στις $20{:}00$ το θερμόμετρο δείχνει $0$° και πέφτει $3$° κάθε ώρα. Η Φανή ρωτά: α) Τι θα δείχνει σε $4$ ώρες; β) Τι έδειχνε $2$ ώρες πριν; Γράψε τις πράξεις με τον ρυθμό $-3$.

\ex \textbf{Το παιχνίδι ερωτήσεων.} Κάθε σωστή απάντηση δίνει $+5$ πόντους και κάθε λάθος $-2$. Ο Ηλίας απαντά σωστά σε $7$ ερωτήσεις και λάθος σε $4$. Ο Στέλιος σωστά σε $3$ και λάθος σε $9$. Πόσους πόντους έχει ο καθένας;

\ex \textbf{Ο ανελκυστήρας.} Ο ανελκυστήρας του ξενοδοχείου κατεβαίνει $2$ ορόφους κάθε φορά που πατάμε «κάτω». Ο Νίκος ξεκινά από τον $6$ο όροφο και πατάει «κάτω» $5$ φορές. Σε ποιον όροφο βρίσκεται;

\ex \textbf{Το νερό του χωριού.} Στη δεξαμενή του χωριού, η γιαγιά Μυρτώ βλέπει ότι το νερό μειώνεται κατά $15$ λίτρα κάθε μέρα (ρυθμός $-15$). Πόσα λίτρα λιγότερα υπάρχουν μετά από $6$ μέρες; Πόσα περισσότερα υπήρχαν πριν από $4$ μέρες;

\ex \textbf{Το μοτίβο της Αγγελικής.} Η Αγγελική γράφει: $(-6)\cdot3=-18$, $(-6)\cdot2=-12$, $(-6)\cdot1=-6$, $(-6)\cdot0=0$. Συνέχισε το μοτίβο με $(-6)\cdot(-1)$ και $(-6)\cdot(-2)$. Τι παρατηρείς;

\zoni{Το σκέφτομαι}{ομόσημοι, ετερόσημοι, πρόσημο γινομένου}

\ex Αν $\alpha\beta>0$, τι ξέρεις για τα πρόσημα των $\alpha$ και $\beta$; Και αν $\alpha\beta<0$;

\ex Πόσοι από τους $5$ παράγοντες ενός γινομένου ακεραίων μπορεί να είναι αρνητικοί, ώστε το γινόμενο να είναι θετικό (όλοι οι παράγοντες μη μηδενικοί);

\ex Ο Στέλιος λέει: «Μείον επί μείον κάνει συν, άρα και μείον συν μείον κάνει συν: $(-3)+(-4)=+7$». Τι μπερδεύει;

\ex Βρες δύο ακέραιους με γινόμενο $-12$ και άθροισμα $1$.

\zoni{Για το σπίτι}{πίνακες και παραδείγματα}

\ex Φτιάξε έναν πίνακα πολλαπλασιασμού με τους ακέραιους από το $-3$ ως το $3$ σε γραμμές και στήλες. Τι παρατηρείς στη γραμμή και στη στήλη του $0$;

\ex Βρες ένα δικό σου παράδειγμα από την καθημερινή ζωή που εξηγεί γιατί το γινόμενο δύο αρνητικών είναι θετικό (π.χ. μια ταινία που παίζει ανάποδα).
